\documentclass[tikz,border=8pt]{standalone}
\usepackage{fontspec}
\usepackage{xeCJK}
\IfFontExistsTF{Noto Sans TC}{\setmainfont{Noto Sans TC}\setCJKmainfont{Noto Sans TC}}{\setmainfont{Microsoft JhengHei}\setCJKmainfont{Microsoft JhengHei}}
\xeCJKsetup{CJKglue={\hskip 0.07em plus 0.02em}, CJKecglue={\hskip 0.10em plus 0.02em}}
\usetikzlibrary{arrows.meta,positioning,calc}
\definecolor{paper}{HTML}{FFFDF8}
\definecolor{navy}{HTML}{0F172A}
\definecolor{muted}{HTML}{334155}
\definecolor{line}{HTML}{CBD5E1}
\definecolor{blue}{HTML}{1D4ED8}
\definecolor{bluebg}{HTML}{DCEBFF}
\definecolor{green}{HTML}{007A5A}
\definecolor{greenbg}{HTML}{DCFCE7}
\definecolor{gold}{HTML}{B85C00}
\definecolor{goldbg}{HTML}{FFE8B5}
\definecolor{red}{HTML}{BA1A1A}
\definecolor{redbg}{HTML}{FFE1E1}
\tikzset{
  box/.style={rounded corners=10pt,line width=1pt,align=center,inner xsep=12pt,inner ysep=8pt,font=\fontsize{9.5}{12.8}\selectfont,text=navy},
  title/.style={font=\bfseries\fontsize{18}{24}\selectfont,text=navy,align=center},
  note/.style={font=\fontsize{10.2}{14}\selectfont,text=muted,align=center}
}
\begin{document}
\begin{tikzpicture}[x=1cm,y=1cm,>=Latex,line cap=round,line join=round]
\path[fill=paper] (0,0) rectangle (16,9.6);
\node[title,text width=14.5cm] at (8,8.78) {讓系統改規則，先限制可修改範圍};
\node[note,text width=13.4cm] at (8,7.92) {self-improving harness 不是放任系統亂改，而是把失敗整理成小範圍修改，再用回歸檢查守住邊界。};

\node[box,fill=bluebg,draw=blue,text width=2.55cm] (fail) at (2.3,5.15) {失敗整理\\同類錯誤歸在一起};
\node[box,fill=goldbg,draw=gold,text width=2.55cm] (edit) at (5.55,5.15) {有限度修改\\只碰允許區域};
\node[box,fill=greenbg,draw=green,text width=2.55cm] (test) at (8.8,5.15) {回歸檢查\\舊功能不能壞};
\node[box,fill=redbg,draw=red,text width=2.55cm] (decide) at (12.05,5.15) {人工判斷\\要不要合併};

\draw[->,draw=navy,line width=1pt] (fail)--(edit);
\draw[->,draw=navy,line width=1pt] (edit)--(test);
\draw[->,draw=navy,line width=1pt] (test)--(decide);
\draw[->,rounded corners=12pt,draw=navy,line width=1pt] (decide.south) -- (12.05,3.15) -- (2.3,3.15) -- (fail.south);

\node[box,fill=white,draw=line,text width=3.2cm] at (4.0,2.05) {不碰權限系統};
\node[box,fill=white,draw=line,text width=3.2cm] at (8.0,2.05) {不改評分標準};
\node[box,fill=white,draw=line,text width=3.2cm] at (12.0,2.05) {不刪錯誤紀錄};

\node[note,text width=13cm] at (8,.86) {能改進，不代表什麼都能改；邊界如果被一起改掉，檢查就失去意義。};
\end{tikzpicture}
\end{document}
